% =====================================================================================
% Draft replacement text for the 3.46 cm (helium) sections of oja_template.tex
% 2026-11-26. Every number is from the gas-cell maps (whole box) unless the text says otherwise.
% Sources: VID/He_cell/stats_He.csv, cell_test_ps/PS_cell_He_NNN.out (ORCD), JALPHA_NOTES §7.1 / §7.1b,
% BSD.ipynb cell 2, memory notes (partial He III ionization, run restart).
% Each block names the lines of oja_template.tex (2026-11-25 version) that it replaces.
% =====================================================================================


% -------------------------------------------------------------------------------------
% (0) NEW paragraph: methods, after l. 252 (end of "Brightness temperature").
%     It needs a label on that subsection:  \subsection{Brightness temperature} \label{sec:tb}
% -------------------------------------------------------------------------------------

The 3.46\,cm signal requires one further step. In the diffuse IGM its spin temperature is set by electron collisions with $x_c \ll 1$, so that $\delta T_\mathrm{b}^\mathrm{He} \propto n_e\,n_\mathrm{He\,II}$ (Appendix~\ref{app:he_density}), and the signal is dominated by dense gas that a $640^3$ grid does not resolve. Evaluated on pixel-averaged fields, the mean signal is 0.24 of its gas-cell value at $z = 5$ and 0.14 at $z = 4$, and even a $10240^3$ grid recovers only 0.82 and 0.75. We therefore compute $\delta T_\mathrm{b}^\mathrm{He}$ for each of the $\approx 2.1\times10^{11}$ gas cells of the \lumina full snapshots, from the cell's own density, temperature and ionization state, with $n_e = n_\mathrm{H\,II} + n_\mathrm{He\,II} + 2n_\mathrm{He\,III}$ and $J_\alpha^\mathrm{He}$ of the $640^3$ pixel that contains the cell, and average it onto the $640^3$ grid, weighted by cell volume. All 3.46\,cm maps and statistics in this paper (maps, VIDs, power spectra and bispectra) use these gas-cell maps, which exist at the redshifts of the full snapshots ($z = 11.9$--3.0). The helium ionization fractions are taken from the $640^3$ renders, which agree with the gas-cell values to $\lesssim 10^{-4}$ per pixel. The light cone (Figure~\ref{fig:lightcone}) has no gas-cell counterpart and uses the $640^3$ fields.


% -------------------------------------------------------------------------------------
% (1) Caption of fig:tmap_helium (l. 616)
% -------------------------------------------------------------------------------------

    \caption{Spatial evolution of the 3.46\,cm brightness temperature $\delta T^\mathrm{He}_\mathrm{b}$ (top panels), from the gas-cell maps, and of the singly ionized helium fraction $x_\mathrm{He\,II}$ (bottom panels), from $z = 8$ to $z = 3$ within a $500\,\mathrm{cMpc}$ \lumina volume.}


% -------------------------------------------------------------------------------------
% (2) Helium maps: replaces l. 643, 646, 657 and 660 (four paragraphs; the VID figure stays between them)
% -------------------------------------------------------------------------------------

Figure~\ref{fig:tmap_helium} shows the spatial evolution of the 3.46\,cm brightness temperature $\delta T_\mathrm{b}^\mathrm{He}$ (upper panels) and of the singly ionized helium fraction $x_\mathrm{He\,II}$ (lower panels) from $z = 8$ to $z = 3$. The signal is weak: its volume average peaks at $0.094\,\mu\mathrm{K}$ at $z \approx 5.5$, about five orders of magnitude below the 21\,cm signal during the EoR. Two factors set this amplitude: the low abundance of $^3$He ($[^3\mathrm{He}/\mathrm{H}] \sim 10^{-5}$) and the weak coupling of the $^3\mathrm{He}^+$ spin temperature to the gas.

Before hydrogen reionization is complete ($z \gtrsim 6$), \ion{He}{2} and free electrons exist only in the ionized regions, because the stellar radiation that ionizes hydrogen also singly ionizes helium. The 3.46\,cm signal therefore fills the \ion{H}{2} regions and is absent from the neutral IGM (the lowest colour level in the upper panels at $z = 8$--6), and the lower panels show the corresponding \ion{He}{2} regions. On large scales the signal is a map of these regions: at $k = 0.1\,h\,\mathrm{cMpc}^{-1}$ its cross-correlation coefficient with $x_\mathrm{He\,II}$ is 0.98--0.99 from $z = 11.9$ to $z = 7$. Inside the regions it follows the cosmic web.

Once hydrogen reionization is complete ($z \lesssim 5.5$), helium is singly ionized almost everywhere and the 3.46\,cm signal covers the whole volume ($z = 5$ panel). It now traces the gas density. Because the spin temperature is set by electron collisions with $x_c \ll 1$, $\delta T_\mathrm{b}^\mathrm{He} \propto x_\mathrm{He\,II}\,x_e\,(1+\delta)^2\,T_\mathrm{K}^{-1/2}$, and the signal correlates with the gas density with $r(k) = 0.90$--$0.99$ and a relative bias of $\approx 2$ (Appendix~\ref{app:he_density}). The \ion{He}{2} Ly$\alpha$ background contributes 7 per cent of the mean signal at $z = 5$ and 18 per cent at $z = 4$. The $(1+\delta)^2$ weighting concentrates the signal in dense gas: gas with $1+\delta > 10$ produces 73 per cent of the mean signal at $z = 5$.

Below $z \approx 5$, quasars doubly ionize helium. The \ion{He}{3} regions appear in the lower panels as dark patches that grow around the quasar hosts and merge, and in the upper panels as regions where the signal drops by several orders of magnitude. The signal does not vanish inside them: dense filaments, where recombinations are faster, keep part of their \ion{He}{2} and remain visible as a network of emission at $z = 3.5$ and $z = 3$. At $z = 3.4$ the mass-weighted \ion{He}{3} fraction is 0.83, against 0.93 by volume. The mean signal therefore falls only from $0.090\,\mu\mathrm{K}$ at $z = 5$ to $0.059\,\mu\mathrm{K}$ at $z = 3.4$ and $0.044\,\mu\mathrm{K}$ at $z = 3$, while $\bar{x}_\mathrm{He\,II}$ falls from 0.9 to below 0.1. Two steps in this evolution come from the simulation set-up rather than from helium reionization (Section~\ref{sec:rt}). At the restart at $z = 4.75$, the uniform UV background ionizes and heats the dense gas and raises the mean signal by 7--9 per cent relative to the trend of the first run. At $z = 3.35$, the onset of self-shielding makes the densest gas partly neutral and lowers the mean signal by 20 per cent, from 0.059 to $0.047\,\mu\mathrm{K}$.


% -------------------------------------------------------------------------------------
% (3) Caption of fig:vid_helium (l. 652)
% -------------------------------------------------------------------------------------

    \caption{Evolution of the VID and its summary statistics for the 3.46\,cm brightness temperature, from the gas-cell maps. Colors indicate redshift from $z = 5$ (dark purple) to $z = 3.4$ (yellow), corresponding to a global ionized fraction $\bar{x}_\mathrm{He\,III}$ increasing from $0.1$ to $0.9$. Top (a): the probability density function. Bottom (b): the redshift evolution of the standard deviation (multiplied by 100), skewness, excess kurtosis and BC from $z = 6$ to $z = 3$, at the redshifts of the full snapshots. The blue dotted line marks $\mathrm{BC} = 5/9$.}


% -------------------------------------------------------------------------------------
% (4) Helium VID: replaces l. 667, 670, 677, 680, 683 (the commented paragraphs between them can go)
% -------------------------------------------------------------------------------------

Figure~\ref{fig:vid_helium} (a) shows the VID of the 3.46\,cm signal from $z = 5$ ($\bar{x}_\mathrm{He\,III} = 0.1$, dark purple) to $z = 3.4$ ($\bar{x}_\mathrm{He\,III} = 0.9$, yellow). Unlike the hydrogen VID, it has no separate emission peak from a neutral phase. At $z = 5$ the distribution is a broad body between $10^{-3}$ and $10^{-1}\,\mu\mathrm{K}$ (median $0.016\,\mu\mathrm{K}$) with a power-law tail that reaches $25\,\mu\mathrm{K}$. Through the $(1+\delta)^2$ weighting, this tail of dense gas carries most of the mean signal.

As the \ion{He}{3} regions grow, an increasing fraction of voxels moves below $\sim 10^{-4}\,\mu\mathrm{K}$, the residual \ion{He}{2} inside the \ion{He}{3} regions, and the median falls by two orders of magnitude, from $0.016\,\mu\mathrm{K}$ at $z = 5$ to $1.9\times10^{-4}\,\mu\mathrm{K}$ at $z = 3.4$. The bright tail does not shrink. Above $\sim 3\,\mu\mathrm{K}$ its probability even rises slightly, because the dense gas that produces it is ionized last. The faint component thus traces the volume of the \ion{He}{3} regions, while the tail follows the \ion{He}{2} left in dense gas.

Figure~\ref{fig:vid_helium} (b) shows the summary statistics from $z = 6$ to $z = 3$. The standard deviation rises from $0.25\,\mu\mathrm{K}$ at $z = 6$ to $0.30\,\mu\mathrm{K}$ at $z = 5$ and $0.39\,\mu\mathrm{K}$ at $z \approx 4$--3.4, although the mean falls from $z \approx 5.5$ on: the signal is carried by fewer, brighter voxels, and the coefficient of variation $\sigma/\langle\delta T_\mathrm{b}^\mathrm{He}\rangle$ grows from 2.8 at $z = 6$ to 6.5 at $z = 3.4$.

The skewness and excess kurtosis are large throughout, far above their hydrogen values: skewness $\approx 11$ and kurtosis $\approx 230$ at $z = 6$--5.5, before helium double ionization begins. They reflect the $(1+\delta)^2$ weighting, which turns the moderately skewed gas density into a strongly one-sided signal. During \ion{He}{2} reionization both rise steadily, with no minimum, to 14.7 and 335 at $z = 4$ and 16.5 and 410 at $z = 3.4$, as the \ion{He}{3} regions remove the diffuse \ion{He}{2} and leave the signal increasingly in dense gas. The 21\,cm statistics have a minimum because the hydrogen field changes from a neutral background with ionized holes to isolated neutral islands. The 3.46\,cm signal has no such inversion, since its bright component remains the dense gas before and after the \ion{He}{3} regions percolate.

The BC rises monotonically, from 0.53 at $z = 6$ to 0.64 at $z = 4$ and 0.70 at $z = 3$, and exceeds the uniform-distribution value of 5/9 from $z \approx 5.2$. For this signal it does not indicate bimodality. With $\mathrm{BC} = (\gamma^2 + 1)/(\kappa + 3)$, where $\gamma$ is the skewness and $\kappa$ the excess kurtosis, its value is set by the heavy one-sided tail ($\gamma \approx 11$--17, $\kappa \approx 230$--410), not by two separated phases.

Two features in panel (b) come from the simulation set-up rather than from helium reionization (Section~\ref{sec:rt}). Between $z = 4.75$ and 4.5, the restart with a uniform UV background ionizes and heats the dense gas, and $\sigma$ rises by 19 per cent. Between $z = 3.37$ and 3.25, the onset of self-shielding makes the densest gas partly neutral, and $\sigma$ falls from 0.38 to $0.31\,\mu\mathrm{K}$ and the kurtosis from 410 to 348.


% -------------------------------------------------------------------------------------
% (5) Caption of fig:ps_helium (l. 750)
% -------------------------------------------------------------------------------------

    \caption{Redshift evolution of the dimensionless power spectra of the doubly ionized helium fraction $\Delta^2_{x_\mathrm{He\,III}}(k)$ (panel a) and of the 3.46\,cm brightness temperature $\Delta^2_\mathrm{He}(k)$ (panel b, gas-cell maps). Colors indicate redshift from $z = 5$ (dark purple) to $z = 3.4$ (yellow), corresponding to a global ionized fraction $\bar{x}_\mathrm{He\,III}$ increasing from $0.1$ to $0.9$. Vertical dashed lines in panel (b) mark $k = 0.1$ (orange) and $0.5\,h\,\mathrm{cMpc}^{-1}$ (purple). Panel (c) shows $\Delta^2_\mathrm{He}$ at these two wavenumbers from $z = 6$ to $z = 3$, at the redshifts of the full snapshots.}


% -------------------------------------------------------------------------------------
% (6) Helium power spectrum: l. 763 and 766 unchanged; replaces l. 768, 771, 780
% -------------------------------------------------------------------------------------

Panel (b) shows the 3.46\,cm power spectrum. At $z = 5$, $\Delta^2_\mathrm{He} \approx 4\times10^{-4}\,\mu\mathrm{K}^2$ at $k = 0.1\,h\,\mathrm{cMpc}^{-1}$ and $\approx 10^{-2}\,\mu\mathrm{K}^2$ at $k = 1\,h\,\mathrm{cMpc}^{-1}$, close to a power law $\Delta^2_\mathrm{He} \propto k^{1.4}$. Unlike the 21\,cm spectrum, it shows no shoulder at the scale of the \ion{He}{3} regions and changes little with redshift. Between $z = 5$ and 3.4 it falls by 30 per cent at $k = 0.1\,h\,\mathrm{cMpc}^{-1}$ and has nearly the same value at $k = 0.5$--$1\,h\,\mathrm{cMpc}^{-1}$, while $\bar{x}_\mathrm{He\,III}$ rises from 0.1 to 0.9. The signal traces the gas density with a relative bias $b_\mathrm{rel} \approx 2$, so that $\Delta^2_\mathrm{He} \approx (b_\mathrm{rel}\langle\delta T_\mathrm{b}^\mathrm{He}\rangle)^2\,\Delta^2_\delta$ (Appendix~\ref{app:he_density}). Over this interval the decline of $\langle\delta T_\mathrm{b}^\mathrm{He}\rangle^2$ (a factor of 0.43) is offset by the growth of the density fluctuations (a factor of $\approx 1.9$). The \ion{He}{3} regions appear only in the comparison with the density: near the midpoint, up to 15 per cent of the power at $k \lesssim 0.15\,h\,\mathrm{cMpc}^{-1}$ is not correlated with the density.

At $k = 0.1\,h\,\mathrm{cMpc}^{-1}$ and $z = 4$ we find $\Delta^2_\mathrm{He} = 3.2\times10^{-4}\,\mu\mathrm{K}^2$. \citet{Spina2025}, who compute the helium spin temperature self-consistently, obtain $\Delta^2_\mathrm{He} \approx 10^{-6}\,\mu\mathrm{K}^2$ at the same scale and redshift. \citet{Basu2026} obtain $\approx 6\times10^{-5}\,\mu\mathrm{K}^2$ at $z = 4.18$ when the spin temperature is set by collisional coupling alone, and $\approx 3\times10^{-2}\,\mu\mathrm{K}^2$ in the fully coupled limit $T_S^\mathrm{He} = T_\mathrm{gas}$. Our value is about five times their collisional estimate, a hundred times below the coupled limit, and about 300 times above \citet{Spina2025}. We attribute most of this difference to resolution. Because the collisional signal scales as $n_e\,n_\mathrm{He\,II}$, it is dominated by dense gas that a uniform grid does not resolve: computed from the $640^3$ pixel-averaged fields instead of the gas cells, our spectrum at $z = 4$ is 50--100 times lower, $(3$--$6)\times10^{-6}\,\mu\mathrm{K}^2$, within a factor of a few of \citet{Spina2025}.

Panel (c) follows $\Delta^2_\mathrm{He}$ at $k = 0.1$ (orange) and $0.5\,h\,\mathrm{cMpc}^{-1}$ (purple). At $k = 0.5\,h\,\mathrm{cMpc}^{-1}$ the power stays between $3.9$ and $5.3\times10^{-3}\,\mu\mathrm{K}^2$ from $z = 6$ to $z = 3.4$. At $k = 0.1\,h\,\mathrm{cMpc}^{-1}$ it halves between $z = 6$ ($8.1\times10^{-4}\,\mu\mathrm{K}^2$) and $z = 5$ ($3.9\times10^{-4}\,\mu\mathrm{K}^2$). This decline marks the end of hydrogen reionization. Before it, \ion{He}{2} exists only in the ionized regions, and the large-scale 3.46\,cm fluctuations follow the pattern of the \ion{H}{2} regions: the relative bias against the density at $k = 0.1\,h\,\mathrm{cMpc}^{-1}$ is 4.9 at $z = 7$ and falls to 2.0 at $z = 5$, once helium is singly ionized everywhere. Below $z = 5$ the large-scale power declines slowly, to $2.6\times10^{-4}\,\mu\mathrm{K}^2$ at $z = 3.5$. The two steps in panel (c) come from the simulation set-up (Section~\ref{sec:rt}). Between $z = 4.75$ and 4.5, the power rises by about 25 per cent at both wavenumbers when the uniform UV background of the restarted run ionizes and heats the dense gas. Between $z = 3.37$ and 3.25, it falls by 31 per cent at $k = 0.5\,h\,\mathrm{cMpc}^{-1}$ when self-shielding makes the densest gas partly neutral. The power does not vanish when helium reionization ends: the \ion{He}{2} left in dense gas keeps $\Delta^2_\mathrm{He} = 2.6\times10^{-4}$ and $2.6\times10^{-3}\,\mu\mathrm{K}^2$ at $k = 0.1$ and $0.5\,h\,\mathrm{cMpc}^{-1}$ at $z = 3$.


% -------------------------------------------------------------------------------------
% (7) Caption of fig:bispec_He_iso (l. 963)
% -------------------------------------------------------------------------------------

    \caption{Redshift evolution of the isosceles bispectrum $\tilde{B}_\mathrm{iso}(k_1, k_2, \theta)$, $k_1 = k_2$, for the singly ionized helium fraction $x_\mathrm{He\,II}$ (left panels) and the 3.46\,cm brightness temperature $\delta T_\mathrm{b}^\mathrm{He}$ (right panels, gas-cell maps), at $k_1 = k_2 = 0.1$, $0.2$, $0.4$ and $1.2\,h\,\mathrm{cMpc}^{-1}$ (top to bottom). Colors indicate redshift from $z = 5$ (dark purple) to $z = 3.4$ (yellow), corresponding to $\bar{x}_\mathrm{He\,III}$ increasing from $0.1$ to $0.9$, with the dashed curve marking the midpoint ($\bar{x}_\mathrm{He\,III} = 0.5$). The left panels use a symmetric logarithmic scale and the right panels a linear scale; the $\delta T_\mathrm{b}^\mathrm{He}$ bispectrum is positive for every configuration. The vertical dashed line marks the equilateral configuration at $\cos\theta = -0.5$.}


% -------------------------------------------------------------------------------------
% (8) Helium isosceles bispectrum: replaces l. 967, 972, 976, 984, 996
% -------------------------------------------------------------------------------------

Figure~\ref{fig:bispec_He_iso} presents the isosceles bispectrum of $x_\mathrm{He\,II}$ (left panels) and $\delta T_\mathrm{b}^\mathrm{He}$ (right panels) at four scales from $k_1 = k_2 = 0.1$ to $1.2\,h\,\mathrm{cMpc}^{-1}$, from $z = 5$ ($\bar{x}_\mathrm{He\,III} = 0.1$) to $z = 3.4$ ($\bar{x}_\mathrm{He\,III} = 0.9$), with the midpoint highlighted as a green dashed curve. In this redshift range helium is singly ionized almost everywhere, so $x_\mathrm{He\,II} \simeq 1 - x_\mathrm{He\,III}$ is the complement of the \ion{He}{3} regions and plays the role that $x_\mathrm{H\,I}$ plays for hydrogen. The \ion{He}{3} regions differ from the \ion{H}{2} regions in two ways: they grow around rare, strongly clustered quasars, and their boundaries are not sharp. At $\bar{x}_\mathrm{He\,III} = 0.1$, three quarters of the volume holds a low level of \ion{He}{3} ($x_\mathrm{He\,III} = 0.01$--0.1), and 35 per cent of all \ion{He}{3} lies in pixels with $x_\mathrm{He\,III} < 0.5$, much of it more than $1\,h^{-1}\,\mathrm{cMpc}$ from the nearest \ion{He}{3} region.

The left panels of Figure~\ref{fig:bispec_He_iso} show a sign reversal in the $x_\mathrm{He\,II}$ bispectrum, as in hydrogen reionization. At early times ($z \sim 5$), the growing \ion{He}{3} regions act as below-average regions in $x_\mathrm{He\,II}$ and drive the bispectrum negative in most configurations. As the midpoint approaches, the contrast diminishes and the negative amplitude decreases. Beyond the midpoint the topology inverts: the remaining \ion{He}{2} forms isolated regions in a \ion{He}{3} background, and the bispectrum turns positive. Before the midpoint the bispectrum has its minimum near the equilateral configuration ($\cos\theta \approx -0.5$), as expected for compact, roughly spherical \ion{He}{3} regions.

After the midpoint the bispectrum turns positive first at the two ends of the angular range. In the stretched limit ($\cos\theta \to 1$) the third leg probes scales smaller than the \ion{He}{3} regions, so the template samples the \ion{He}{2} between them, an above-average structure once \ion{He}{3} fills most of the volume. In the squeezed limit ($\cos\theta \to -1$) the positive signal shows that small-scale \ion{He}{2} fluctuations are stronger where the large-scale \ion{He}{2} fraction is high, that is, in the regions the quasar radiation has not yet reached.

As in the hydrogen analysis, the sign change of the $x_\mathrm{He\,II}$ bispectrum measures a characteristic size of the \ion{He}{3} regions. At $k_1 = k_2 = 1.2\,h\,\mathrm{cMpc}^{-1}$ the sign flip is already present at $z = 5$, near $\cos\theta \approx 0.85$. Section~\ref{sec:bubble} derives these sizes systematically and compares them with the MFP.

The 3.46\,cm bispectrum (right panels) behaves differently. It is positive for every configuration, scale and redshift, and changes much less during helium reionization than that of $x_\mathrm{He\,II}$, which changes sign: in the stretched limit at $k_1 = k_2 = 0.4\,h\,\mathrm{cMpc}^{-1}$ it rises from 15 at $z = 5$ to 19 at $z = 3.4$. Its amplitude grows with $k$, from $\tilde{B} \approx 0.3$--4 at $k_1 = k_2 = 0.1\,h\,\mathrm{cMpc}^{-1}$ to 10--50 at $1.2\,h\,\mathrm{cMpc}^{-1}$, and at each scale it is smallest near the equilateral configuration and largest in the stretched limit. This is the bispectrum of a density tracer. After hydrogen reionization the 3.46\,cm signal follows the gas density ($r = 0.90$--0.99) with a quadratic response, so its bispectrum is set by the non-linear clustering of the gas, enhanced by this response, and is positive. The \ion{He}{3} regions, which drive the sign change of the $x_\mathrm{He\,II}$ bispectrum, modulate the signal only weakly.


% -------------------------------------------------------------------------------------
% (9) Caption of fig:non_iso_He (l. 1014)
% -------------------------------------------------------------------------------------

    \caption{Redshift evolution of the non-isosceles bispectrum $\tilde{B}_\mathrm{non\text{-}iso}(k_1, k_2, \theta)$ for the singly ionized helium fraction $x_\mathrm{He\,II}$ (left panel) and the 3.46\,cm brightness temperature $\delta T_\mathrm{b}^\mathrm{He}$ (right panel, gas-cell maps), evaluated at the fixed configuration $k_1 = k_2/2 = 0.3\,h\,\mathrm{cMpc}^{-1}$. Colors and dashed curve follow the same convention as Figure~\ref{fig:bispec_He_iso}. The right panel uses a linear scale. The vertical dashed line marks the configuration at $\cos\theta = -0.5$.}


% -------------------------------------------------------------------------------------
% (10) Helium non-isosceles bispectrum: l. 1018 unchanged; replaces l. 1023, 1027, 1030
% -------------------------------------------------------------------------------------

Figure~\ref{fig:non_iso_He} shows the redshift evolution of the non-isosceles bispectrum for $x_\mathrm{He\,II}$ (left) and $\delta T_\mathrm{b}^\mathrm{He}$ (right). The $x_\mathrm{He\,II}$ bispectrum shows the same inversion as in hydrogen (Figure~\ref{fig:non_iso_H}), from negative to positive as helium reionization proceeds, but its folded limit ($\cos\theta \to -1$) turns positive early, at $\bar{x}_\mathrm{He\,III} \approx 0.2$--0.3, well before the midpoint, while the intermediate configurations stay negative until $\bar{x}_\mathrm{He\,III} \approx 0.5$--0.7. The early positive cross-scale coupling between $r_1$ and $r_2$ reflects the strong clustering of the quasar hosts, around which the \ion{He}{3} regions and the surrounding \ion{He}{2} structure are concentrated. Toward the stretched limit ($\cos\theta \to 1$) the bispectrum turns positive once the third leg probes scales below the size of the \ion{He}{3} regions, as in the isosceles case.

As in the hydrogen non-isosceles case, the intermediate configurations show an asymmetric trough near $\cos\theta \approx -0.5$, shifted by the broken leg symmetry, which couples the unequal base scales $r_1 \approx 10\,h^{-1}\,\mathrm{cMpc}$ and $r_2 \approx 5\,h^{-1}\,\mathrm{cMpc}$. Its location is fixed by the triangle geometry, but its depth evolves: the early \ion{He}{3} regions deepen the negative signal, and after the midpoint the trough rises to positive values while remaining the minimum.

The 3.46\,cm non-isosceles bispectrum is positive for every configuration and redshift, as in the isosceles case, and rises slowly from $z = 5$ to 3.4 (from 16 to 21 in the stretched limit, from 6 to 8 at its minimum near $\cos\theta \approx -0.5$). Since the signal traces the gas density with a quadratic response, its bispectrum reflects the non-linear clustering of the gas rather than the topology of the \ion{He}{3} regions. This differs from hydrogen, where the 21\,cm bispectrum follows the sign change of the $x_\mathrm{H\,I}$ bispectrum: the 21\,cm signal is set by the neutral fraction, which varies by order unity between ionized and neutral regions, whereas after hydrogen reionization the 3.46\,cm signal varies mainly with the density.


% -------------------------------------------------------------------------------------
% (11) Helium bubble statistics: replaces l. 1109, 1114, 1118, 1122, 1125
% -------------------------------------------------------------------------------------

Figure~\ref{fig:bubble_size} (b) presents the MFP \ion{He}{3} bubble sizes from $z \approx 5.2$ to $z \approx 3.1$, computed as for hydrogen with ionized cells defined by $x_\mathrm{He\,III} > 0.5$. As in the hydrogen analysis, the mean lies above the median (blue dashed and solid curves), a skewed distribution dominated by rare, large quasar-driven regions. Overlaid are the estimates from the power spectrum peak (purple circles) and the bispectrum sign flip (green diamonds).

The median MFP radius grows from 4.2\,cMpc at $z = 5.0$ to 13.6\,cMpc at $z = 4.0$ and 106\,cMpc at $z = 3.44$, where the 84th percentile reaches 292\,cMpc. The broad percentile range shows that small and very large \ion{He}{3} regions coexist throughout, as expected from the highly biased quasar population that drives helium reionization.

Unlike in hydrogen reionization, the power spectrum estimates exceed the MFP scales from the outset and grow far more slowly. They start at $\approx 12$\,cMpc at $x_\mathrm{He\,II} = 0.9$, about three times the MFP median, and grow only to $\approx 28$\,cMpc by $x_\mathrm{He\,II} = 0.1$, while the MFP median grows by more than an order of magnitude. Because quasars are rare and strongly clustered within massive overdensities, the \ion{He}{3} regions around neighbouring quasars overlap, so the power spectrum peak likely traces the large-scale clustering of these quasar environments, whereas the MFP measures the much smaller distances to individual \ion{He}{3} boundaries. As in hydrogen, the power spectrum estimate then saturates while the MFP sizes keep growing, and it falls below the median for $x_\mathrm{He\,II} \lesssim 0.35$ as the \ion{He}{3} regions percolate.

For the bispectrum sign flip we take the nearly stretched isosceles configuration ($k_3 \approx 2k$) and find the $k$ at which the $x_\mathrm{He\,II}$ bispectrum changes sign, $R_\mathrm{ion} = \pi/k_{3,t}$; where it changes sign twice, we take the geometric mean of the two crossings. At $\bar{x}_\mathrm{He\,III} = 0.10$ and 0.15 ($z = 5.0$ and 4.75) this gives 2.8 and 2.9\,cMpc, about 0.6--0.7 times the MFP median (4.1 and 5.0\,cMpc). At $\bar{x}_\mathrm{He\,III} = 0.21$ and 0.33 ($z = 4.5$ and 4.25) it gives 8.8 and 15\,cMpc, close to or above the MFP mean. The low early values come from the partial ionization described in Section~\ref{sec:bispectrum}: the continuous $x_\mathrm{He\,II}$ field includes the widespread low-level \ion{He}{3} outside the \ion{He}{3} regions, which the MFP, with its threshold $x_\mathrm{He\,III} > 0.5$, does not count. Applied to the binary field defined by the same threshold, the method gives 6.4 and 7.7\,cMpc, close to the MFP mean (6.1 and 7.5\,cMpc). From $\bar{x}_\mathrm{He\,III} \approx 0.2$ the two fields give the same result. As in hydrogen reionization, the sign flip no longer yields a clean \ion{He}{3} scale once percolation inverts the topology.

For comparable global ionized fractions, the \ion{He}{3} regions are systematically larger than their \ion{H}{2} counterparts: at an ionized fraction of 0.5 the MFP median is $\approx 14$\,cMpc for \ion{He}{3} and $\approx 6$\,cMpc for \ion{H}{2}. This difference stems from the source populations. Hydrogen reionization is powered by faint, abundant stellar sources, whereas helium reionization is driven by rare, luminous quasars whose hard spectra and long photon mean free paths produce larger, more strongly clustered ionized regions. The distinction fades near the end of reionization, as the regions percolate into box-scale networks.


% -------------------------------------------------------------------------------------
% (12) Sub-volume VID, helium: replaces l. 1191. Caption of fig:subbox (l. 1163) is correct as it stands.
% -------------------------------------------------------------------------------------

Figure~\ref{fig:subbox} (b) shows the corresponding VID of the 3.46\,cm signal at $z = 5$, where $\bar{x}_\mathrm{He\,III} = 0.1$ as in the hydrogen case. Every 100\,cMpc sub-volume reproduces the bulk of the distribution and its bright tail, which carries most of the mean signal. The sub-volumes scatter only below $\sim 10^{-3}\,\mu\mathrm{K}$, by up to two orders of magnitude in probability density, in the faint voxels of the low-density IGM and of the first \ion{He}{3} regions. How many such voxels a sub-volume contains depends on whether it holds one of the rare quasars that start helium double ionization. Quasars are far rarer and more strongly clustered than star-forming galaxies, and a 100\,cMpc volume does not sample them representatively, whereas the 500\,cMpc \lumina volume does.
